% Owner: environment / benchmark (this chapter)
% Do not put 120+ or unofficial run scores here.
% Cross-refs: \ref{sec:intro}, \ref{sec:related}, \ref{sec:exp},
%   \ref{tab:mgpp-compare}, \ref{app:mgpp-apps}, \ref{app:mgpp-tools}
\section{MobileGym++}
\label{sec:mobilegym++}
\label{sec:method}

\subsection{A gap in controllable commercial applications}
\label{sec:mgpp-env}

Mobile agents are moving from system utilities and open-source clients into a larger commercial application world. Travel booking, online shopping, personal finance, and local services occupy most of a user's daily time, and they are the settings agents must eventually enter. Scaling evaluation and training into those settings, however, hits a structural bottleneck on each of the two paths now available.

One path runs real applications on physical devices. It buys interface and network fidelity at the cost of device state that is hard to reset or parallelize. Many commercial apps refuse to load once they detect an automation environment, so collection is pinned to particular devices and accounts, and an order, a payment, or a sent message during an experiment is an irreversible write into the real world. The other path runs reproducible applications in a controllable environment. It buys reset and a deterministic checker at the cost of an inclusion rule that favors offline, login-free software---mostly system utilities and open-source clients. AndroidWorld's rewards come from system state read through debugging interfaces: they can confirm a calendar event or a contact change, but not the amount paid or the recipient on a third-party shopping order~\citep{rawles2025androidworld}. AppWorld supplies a controllable world and state-based programmatic evaluation, but its action space is function calls and there is no operable graphical interface~\citep{trivedi2024appworld}. MobileWorld adds cross-application work and external tools, yet many items are written so that they are hard to finish without those tools~\citep{kong2025mobileworld}. In earlier controllable suites, the terminals a program can check concentrate on system settings and local files. Completing an order, posting an expense, or changing an itinerary still has nowhere stable to land.

Judgment choices shrink the usable range further. A model judge introduces unstable false positives. Once an online site is taken offline, its native verifier no longer applies and comparison falls back to matching a reference trajectory, whose gap from other valid paths is hard to close. For a setting that is meant to serve both academic and industrial use, the check has to land on the application's own business objects---and that is possible only when the application itself is fully controllable.

Neither path, then, carries scalable evaluation and training through a complete commercial operation. The device path cannot reset state or undo side effects; the controllable path is narrow in application mix and shallow in what the terminal proxy can see. Covering commercial work depends on constructing controllable mock commercial applications.

MobileGym++ is the response: a browser-hosted mobile environment that reproduces closed commercial loops. It ships 64 installable applications, 35 of them newly built commercial clients, covering travel, shopping, finance, content, social, office, and local services. Authenticity, diversity, and stability are the acceptance criteria. Authenticity is the closed loop---goods and orders, expenses and budgets, itineraries and changes exist as objects in each app's own store. Diversity is the seven domains and 64 applications. Stability is an initial state that can be exported and restored as a whole, together with a programmatic checker that reads only business state. The runtime is MobileGym~\citep{wu2026mobilegym}. Serialization, reset, inter-app navigation, and programmatic checking are inherited. What this paper expands is the content those mechanisms carry, and one new calling channel.

\subsection{Platform and expansions}
\label{sec:mgpp-inherit}

The predecessor already supplies a complete controllable runtime. The device runs inside a browser tab; each app keeps a structured store and renders its UI from that store; the clipboard, system broadcasts, and inter-app navigation already connect apps; the full device state can be exported as one record and restored from it; many instances can run at once without interfering. On the agent side, observations remain screen pixels. There is no accessibility tree, and the agent cannot read the business store. Actions remain discrete taps, swipes, typing, and navigation; at most one action executes per step.

On that same runtime we make two expansions. The first spreads reproducible commercial clients to a scale that covers domains. Newly built apps reproduce a business loop rather than a skin. Each client is built from a real product: the object model and state transitions are fixed first, then the business functions and pages. Goods and orders, expenses and budgets, itineraries and changes, conversations and messages exist as objects in the app's own store. The UI renders them; the checker reads them. Before admission, each newly built app must pass a page-reachability check and a terminal-state checkability check. The second expansion opens, inside the apps, a set of function calls with the same rights as the GUI; we return to it in Section~\ref{sec:mgpp-tools}.

\subsection{Programmatic checks on business state}
\label{sec:mgpp-check}

A task binds an initial state. At the end we take a diff of the full device state and ask whether the target business fields were written as required. The checker depends on neither a model judge nor a single reference trajectory: different routes that accomplish the same write are equivalent. Programmatic state checking is already shared practice in controllable environments. AppWorld contrasts state-based evaluation with the brittleness of reference matching~\citep{trivedi2024appworld}; AndroidWorld notes that system state outlasts matching surface UI changes~\citep{rawles2025androidworld}; MobileGym states that checks are deterministic, fine-grained, and independent of a vision-language judge~\citep{wu2026mobilegym}. The difference here is not the mechanism but the object being read. The checker reads the app's own business store, so a commercial result can enter evaluation. Treating a system setting or a widget tree as a terminal proxy cannot cover an operation whose landing point is a business object.

Checks are written as relations, not as constants. An expense amount must equal the payment that just completed; a calendar title must contain the selected train number. Constants appear only in the initial state. Outside the target fields, a task declares which state paths may change; remaining diffs are recorded as unexpected mutations and reported separately from success. A task counts as cross-application if and only if the checker reads a business write in at least two apps. Querying and then paraphrasing, or sending a notification into a second app, does not count.

\subsection{Scale, depth, and coverage}
\label{sec:mgpp-scale}

The environment ships 64 installable applications, 35 of them newly built commercial clients, organized by the seven domains above. Each domain carries a different object type and a different terminal form---placing an order, posting a ledger entry, creating an itinerary, sending a message, adding a favorite---so the same checker template cannot cover them all.

Count alone is not enough; each app must be deep enough to carry a multi-step task. On the newly built clients, the median navigation declaration is 37 routes, 175 interactive actions, and 88 transitions. As a depth reference, the thickest application in MobileGym has 100 routes and 265 actions, and its deepest terminal takes 7 steps to reach. Median interaction depth from the root in this environment is 3: the shortest GUI path of a typical task already crosses several layers, and the step-count lower bound comes from page structure rather than from the length of the instruction. Table~\ref{tab:mgpp-compare} summarizes how this environment differs from other controllable gyms in what a program is allowed to read. Per-app statistics are in Appendix~\ref{app:mgpp-apps}.

\begin{table}[t]
\caption{Controllable environments. The checker column is what a program is allowed to read; the tool column is whether a call writes the same store as the GUI. Real-device gyms are omitted: they cover vendor apps and live networks, not a resettable commercial store.}
\label{tab:mgpp-compare}
\begin{center}
\scriptsize
\setlength{\tabcolsep}{3pt}
\begin{tabular}{lcccc}
\hline
 & Apps & Checker reads & Cross-app & Tools vs.\ GUI \\
\hline
AndroidWorld & system / default & system state & limited & --- \\
AppWorld & APIs, no phone GUI & application DBs & yes & API only \\
MobileWorld & longer vendor-style & mixed & yes & external MCP \\
MobileGym & mock daily / system & device state diff & OS-level & none \\
MobileGym++ & 64 mock commercial & app business store & dual writes & same store \\
\hline
\end{tabular}
\end{center}
\end{table}

\subsection{In-app tools}
\label{sec:mgpp-tools}

Agents already call functions. Existing evaluation splits that fact across two settings. Graphical environments almost only score taps and typing; interface-only environments have nothing to tap. The behavior that then has nowhere to be tested is the choice itself: when to use the GUI, when to call an interface, and how to return to the screen to confirm. As an extension for future evaluation, this environment places both channels on the same business store, in one runtime. GUI-only research remains valid; no official task path can be finished only through tools.

Approximately 52 applications expose a function interface under their own namespace. The agent may list and call the foreground app's functions. A call occupies the same action slot as a GUI action; at most one of the two executes per step. Those functions act on the app's own store and go through the same business functions as the GUI, so they are a second channel in one world rather than an external service that bypasses the screen. The operations they cover are a strict subset of what the GUI can already finish. Payment confirmation, password entry, and contact picking remain GUI-only. Prefill calls produce a draft and no business write; saving still requires the screen.

A typical mixed solution looks like this. The agent calls the bookkeeping app to read the currently available budget categories, switches to the shopping app, picks an item and completes payment on screen, then calls the bookkeeping app again to write that expense into a named ledger. Payment cannot be finished by a call, because the PIN pad exists only in the GUI. Bookkeeping admits both routes: a save on screen and a function call produce the same record. The checker reads the same fields in either case, and therefore does not distinguish whether the expense was tapped or invoked. Every official task keeps a GUI-only path. A GUI-only solution and a mixed solution are compared on the same task, the same initial state, and the same checker. Protocol and the tool catalog are in Appendix~\ref{app:mgpp-tools}.

\subsection{Benchmark}
\label{sec:mgpp-bench}

\paragraph{Task design.}
Single-app GUI operation is already scored thoroughly. From AndroidWorld to MobileWorld, the typical task is a sequence of interface actions inside one application, and the terminal lands on system-level state or a surface attribute of the app---editing a contact, toggling a setting, creating a calendar event~\citep{rawles2025androidworld,kong2025mobileworld}. Along that axis, existing coverage reduces to an in-app interface chain whose terminal is proxied by system state or a widget tree. That axis is densely and maturely scored. It does not reach two further places.

The first is a business write as the terminal. Completing an order, posting an expense, or changing an itinerary does not land in a system setting or on a widget tree; it lands in the app's own business store. On the real-device path, that store cannot be read in full by an external program, so checking falls back to screenshot matching or a model judge. On the controllable path, earlier environments stop at system utilities and open-source clients, and the business store itself is absent. Operations whose landing point is a business object therefore still lack a programmatic terminal.

The second is a cross-application value dependence. In a single-app task, everything needed to finish is usually written in the instruction or already visible on the current screen. In a commercial process, the number a write needs often comes from another app's live state: the amount actually paid after checkout must be posted to a ledger; the selected train must be copied into a calendar. The agent must read the value in one app and write it into the correct field in another. Cross-application items in existing suites mostly stop at querying and then paraphrasing, or at sending a notification; the second app does not produce a business write that depends on a fact from the first, and no checker can ask whether the two apps' fields line up.

MobileGym++ Bench sits on the environment above and occupies exactly those two extensions. Every task ends in a business write that actually occurs. Tasks that only query, only open a page, or only restate a result in text do not enter the official set. Realism comes from the process, not from the wording of the prompt: a task is long because it must pass through selection, confirmation, payment, and write-back, not because unrelated steps were stacked. The instruction is written as a user would speak it; it does not name widget labels, in-app namespace identifiers, or checker fields.

In a cross-application task the needed number is absent from the instruction and present in another app's live store. The instruction gives a pointer---``the amount actually paid'', ``the largest charge'', ``this train''---and the checker asks whether the two apps' fields line up, not whether a constant written in advance appears. The relation moves with the environment state and cannot be solved by memorizing constants in the prompt.

A GUI-only solution is a first-class citizen. Every task keeps at least one GUI-only path; a mixed solution is another path in the same world, reported under the same checker.

Difficulty does not come from a longer prompt or from stacked unrelated steps. It comes from the structural cost of finishing a write. Median shortest GUI path in the environment is 3; deeper tasks enter further layers, and that lower bound is set by page structure. The same business action can look finished on screen and land in different store fields: ledger category, order ownership, and message channel are distinct. The checker reads the field; stopping on a result page is not evidence of completion. A cross-application task additionally requires the agent to read the needed value in one app and write it in another, and both stores are read by the checker.

\paragraph{Task specification.}
Each task has three published parts: a natural-language user goal, a bound initial state, and a programmatic check on the terminal state. The checker reads the before/after diff of business state. It does not read the trace, does not read screenshots, and does not accept an agent's claim of completion. Paths that the task allows to change are declared; other diffs are reported as unexpected mutations, separately from success.

\paragraph{Scale and composition.}
Tasks are organized by business domain, not by application, and cover the same seven domains as the environment. The design size is approximately 200 tasks, of which approximately 116 are single-app and 84 are cross-application under the write-based definition above. Each frequently used app carries several tasks, and a few high-frequency apps carry more, so that evaluation is not dominated by a handful of apps and no app is covered by a single item. Only tasks that meet the checker-level definition count as cross-application---a business write in at least two apps. An app that appears only as a notification channel is not counted toward its domain's cross-application coverage.

\paragraph{Evaluation protocol.}
The primary metric is task success rate. Step counts are reported on successful tasks as an efficiency reference; failed trajectories often stop before the step cap, and those steps are excluded from the efficiency statistic. The main table reports two success rates on the same tasks, the same initial state, and the same checker: GUI-only and mixed. They have equal status. The release includes the task definitions, the initial-state records, and the read paths the checker uses. Runtime version and the time base used in evaluation are locked in the protocol so that runs remain comparable.
